
\medskip
\phantomsection\label{ctr:sga3-VIB-1.4.3}%
\noindent\textbf{Counter\SGARefLink{ctr:sga3-VIB-1.4.3}{example 1.4.3}.---}
\emph{Let \(k\) be a field of characteristic zero, let \(G\) be
the constant \(k\)-group \(\mathbb Z\), and let \(H\) be the
additive \(k\)-group \(\mathbb G_{a,k}\). Let \(u:G\to H\) be a
morphism of \(k\)-groups. If \(u\ne0\), then \(\ker u=0\), but
\(u\) is not a closed immersion (cf.\ \hyperref[ctr:sga3-VIB-1.2.2]{1.2.2}).}

We shall use below the following two results, which should have
appeared in Expos\'e~VI\(_A\).

\medskip
\phantomsection\label{lem:sga3-VIB-1.5}%
\noindent\textbf{Lemma 1.5.---}
\emph{Let \(k\) be a field and let \(G\) be a \(k\)-group locally
of finite type. Then:
\begin{enumerate}
\item[\textup{\SGAInlineRefTarget{lem:sga3-VIB-1.5:item:i}(i)}] all the irreducible components of \(G\) have
  the same dimension;
\item[\textup{\SGAInlineRefTarget{lem:sga3-VIB-1.5:item:ii}(ii)}] for every \(g\in G\),
  \[
    \dim_gG=\dim G.
  \]
\end{enumerate}}

\footnote{\label{note:sga3-VIB-17}Editor's note: the implication
\textup{(ii)} \(\Rightarrow\) \textup{(i)} is a general fact
(cf.\ EGA~0\(_{\mathrm{IV}}\), 14.1.6), while
\textup{(i)} \(\Rightarrow\) \textup{(ii)} follows from the fact
that if \(X\) is an irreducible scheme of finite type over a field,
then \(\dim X=\dim U\) for every nonempty open subset \(U\) of
\(X\). The essential point here is therefore assertion \textup{\SGARefLink{lem:sga3-VIB-1.5:item:i}{(i)}},
which was already established in an addition to
Expos\'e~VI\(_A\), \SGARefLink{cor:sga3-VIA-2.4.1}{2.4.1}. We have modified the statement and proof
accordingly.}
Assertion \textup{\SGARefLink{lem:sga3-VIB-1.5:item:i}{(i)}} has already been proved in
Expos\'e~VI\(_A\), \SGARefLink{cor:sga3-VIA-2.4.1}{2.4.1}, and \textup{\SGARefLink{lem:sga3-VIB-1.5:item:ii}{(ii)}} follows. Indeed, let
\(g\in G\), and let \(C\) be the connected component of \(G\)
containing \(g\). By definition
(EGA~0\(_{\mathrm{IV}}\), 14.1.2), \(\dim_gG\) is the infimum of
the integers \(\dim U\), where \(U\) ranges over the open
neighborhoods of \(g\). Hence
\[
  \dim_gG=\dim U_0
\]
for some \(U_0\), which we may suppose is contained in \(C\)
(because \(\dim V\leq\dim U\) when \(V\subset U\)). Since \(C\) is
irreducible and of finite type over \(k\)
(Expos\'e~VI\(_A\), \SGARefLink{cor:sga3-VIA-2.4.1}{2.4.1}),
\[
  \dim U_0
    =\dim C
    =\operatorname{trdeg}_k\kappa(\xi),
\]
where \(\xi\) is the generic point of \(C\)
(EGA~IV\(_2\), 5.2.1). Therefore
\[
  \dim_gG=\dim C=\dim G.
\]

\medskip
\phantomsection\label{prop:sga3-VIB-1.6}%
\noindent\textbf{Proposition 1.6.---}%
\footnote{\label{note:sga3-VIB-18}Editor's note: in the statement,
we have replaced ``along the identity section'' by ``at every point
of the identity section.'' At the end of the proof, we have also
made explicit the cited results from EGA~IV\(_4\).}
\emph{Let \(S\) be a scheme of characteristic zero, and let \(G\)
be an \(S\)-group scheme locally of finite presentation over \(S\)
at every point of the identity section \(\epsilon(S)\). The scheme
\(G\) is smooth over \(S\) at every point of the identity section
if and only if the \(\mathcal O_S\)-module}
\[
  \omega_{G/S}
    =\epsilon^*\bigl(\Omega^1_{G/S}\bigr),
\]
\emph{called the conormal module of the identity section of \(G\),
is locally free.}

Recall that a scheme \(S\) is said to be of characteristic zero if
the field \(\kappa(x)\) has characteristic zero for every closed
point \(x\in S\).

Recall also (\SGARefLink{sga3:ii:subsection:4-11}{Expos\'e~II, 4.11}) that, if
\(\pi:G\to S\) is the structural morphism, then
\[
  \Omega^1_{G/S}=\pi^*(\omega_{G/S}).
\]
Thus it is equivalent to say that the \(\mathcal O_S\)-module
\(\omega_{G/S}\) is locally free or that the
\(\mathcal O_G\)-module \(\Omega^1_{G/S}\) is locally free.

If there is an open neighborhood \(U\) of \(\epsilon(S)\) that is
smooth over \(S\), then by EGA~IV\(_4\), 17.2.3,
\(\Omega^1_{U/S}\) is locally free of finite type, and so is
\[
  \omega_{G/S}=\epsilon^*(\Omega^1_{U/S}).
\]

Conversely, if \(\omega_{G/S}\) is locally free, then so is
\[
  \Omega^1_{G/S}=\pi^*(\omega_{G/S}).
\]
Since \(S\) has characteristic zero, the Jacobian criterion
(EGA~IV\(_4\), 16.12.2) shows that \(G\) is differentially smooth
over \(S\). It follows from EGA~IV\(_4\), 17.12.5, that \(G\) is
smooth over \(S\) at every point of the identity section.

\medskip
\phantomsection\label{cor:sga3-VIB-1.6.1}%
\noindent\textbf{Corollary 1.6.1 (Cartier).---}
\emph{Let \(k\) be a field of characteristic zero. Every
\(k\)-group locally of finite type over \(k\) is smooth over
\(k\).}

Indeed, the \(k\)-module \(\omega_{G/k}\) is then plainly free.
By \SGARefLink{prop:sga3-VIB-1.6}{1.6}, \(G\) is smooth over \(k\) at the identity point \(e\), and
hence smooth over \(k\) by \hyperref[cor:sga3-VIB-1.3.1]{1.3.1}.
